\section{Antecedentes: acceso, infraestructura y continuidad}
La experiencia venezolana con ATLAS Open Data muestra que los datos abiertos pueden utilizarse en trabajos de pregrado y posgrado con herramientas de análisis y acompañamiento académico \cite{atlas}. La existencia de un conjunto de datos es una condición de partida; una actividad formativa necesita además ejemplos ejecutables, documentación y una pregunta adecuada a su duración. La accesibilidad debe evaluarse en esa combinación.

Los materiales de contexto aportados por el autor identifican dos antecedentes técnicos: la contribución CHEP 2018 sobre Open Key y el despliegue multiusuario propuesto en CHEP 2021 \cite{contexto,chep}. Su relación con el presente trabajo es la búsqueda de un punto de partida utilizable para estudiantes. Un entorno transportable, un servicio compartido y un espacio en la nube resuelven problemas parcialmente distintos. No se plantea una sustitución universal de los recursos locales por servicios externos.

El documento de estrategia de 2024 propone aprovechar recursos y capacidades compartidos para investigación y formación \cite{whitepaper}. El presente manuscrito se concentra en una escala operativa: un repositorio de ejercicios, un entorno descrito explícitamente y una secuencia de análisis. La contribución consiste en relacionar esas piezas con el proceso docente y con un protocolo de verificación. La portabilidad se considera una propiedad que debe probarse, no una consecuencia automática de usar Git.

\subsection{De ROOT en la web a la formación comunitaria en software}
Una presentación de Sánchez Pineda en el ROOT Workshop de septiembre de 2015 ya relacionaba interfaces web para analizar datos de ATLAS con tutoriales de ROOT y actividades de CEVALE2VE \cite{root2015}. Este antecedente sitúa el acceso mediante navegador dentro de una trayectoria anterior a Codespaces. Las herramientas y los despliegues han cambiado, pero permanece una pregunta docente: cómo acercar el análisis de física a estudiantes con recursos y experiencia heterogéneos.

El artículo colectivo \emph{Software Training in HEP}, del que Sánchez Pineda es coautor, organiza la formación en competencias generales de software, herramientas específicas de física y técnicas avanzadas \cite{hsftraining}. Esa distinción permite situar Git y la integración continua como competencias transferibles, junto con herramientas disciplinares como ROOT. El presente trabajo desarrolla una aplicación docente de esa combinación; no propone que el uso de una plataforma concreta constituya por sí mismo una competencia científica.

\subsection{ESCAPE y los entornos virtuales de investigación}
ESCAPE vincula infraestructuras de física de partículas y astrofísica con servicios de ciencia abierta \cite{escape}. Su directorio de herramientas reúne accesos a JupyterHub, entornos Docker, documentación de Rucio, repositorios de software, tutoriales y recursos de ATLAS Open Data \cite{escapetools}. El Virtual Research Environment documentado por VRE Hub busca facilitar flujos de análisis completos y la preservación de resultados conforme a principios FAIR \cite{vrehub}.

Estas fuentes permiten comparar dos escalas de trabajo. Un entorno virtual de investigación articula servicios y recursos de una infraestructura compartida; el repositorio docente aquí descrito delimita un ejercicio y sus dependencias para una sesión de formación. Comparten la necesidad de documentar el entorno, pero no son sistemas equivalentes. No se atribuye a los talleres de este estudio el uso de los servicios de ESCAPE sin registros que lo documenten. La contribución de esta comparación es identificar requisitos reutilizables: acceso a datos, descripción del software y conservación del procedimiento.
